\documentclass[11pt]{article}

% ---------- Packages ----------
\usepackage[margin=1in]{geometry}
\usepackage{hyperref}
\usepackage{longtable}
\usepackage{enumitem}
\usepackage{courier}

\setlength{\parskip}{6pt}
\setlength{\parindent}{0pt}

\title{\textbf{Novel Knowledge System (NKS)}\\Comprehensive CLI Documentation}
\author{}
\date{March 2026}

\begin{document}

\maketitle

\section{Introduction}

The \textbf{Novel Knowledge System (NKS)} is a CLI-based tool designed to manage structured data for novel writing. It integrates:

\begin{itemize}
    \item A relational SQL database for persistence
    \item A graph-based relationship system for character links
    \item A command-line interface optimized for real-time workflows
\end{itemize}

This document provides a complete reference of the supported commands, including syntax, explanations, and examples.

\section{Command Format}

All commands follow a consistent structure:

\begin{verbatim}
novel <entity> <action> <arguments> [options]
\end{verbatim}

When running the application without arguments, NKS starts an interactive shell with a \texttt{novel>} prompt. In that shell, commands may be entered either with the \texttt{novel} prefix or without it, and the session remains active until \texttt{exit} is entered.

Placeholders are represented as:

\begin{itemize}
    \item \texttt{<novel>} -- Novel title
    \item \texttt{<id>} -- Numeric identifier
    \item \texttt{<name>} -- Entity name
    \item \texttt{<text>} -- Free-form content
\end{itemize}

\section{1. Initialization and Configuration}

\subsection{Interactive Shell}

\textbf{Syntax}
\begin{verbatim}
python3 project/main.py
\end{verbatim}

\textbf{Description}

Starts an iterative CLI session powered by \texttt{prompt\_toolkit}. The shell provides tab-completion for command phrases and common stored values such as novel titles, section names, and character names.

\textbf{Example}
\begin{verbatim}
novel> help
novel> create "The Great Swordman"
novel> review report "The Great Swordman"
novel> exit
\end{verbatim}

\subsection{novel init}

\textbf{Syntax}
\begin{verbatim}
novel init [--db <database_url>]
\end{verbatim}

\textbf{Description}

Initializes the system by creating the database schema and writing a local configuration file. A database URL must be supplied by \texttt{--db}, \texttt{NKS\_DATABASE\_URL}, or the gitignored local config file.
\newline
\newline
For privacy, connection credentials should be supplied through an explicit \texttt{--db} value, the \texttt{NKS\_DATABASE\_URL} environment variable, or the gitignored local config file.

\textbf{Example}
\begin{verbatim}
NKS_DATABASE_URL="postgresql+psycopg://<user>:<password>@localhost:5432/novel_notes" novel init
\end{verbatim}

\subsection{novel migrate}

\textbf{Syntax}
\begin{verbatim}
novel migrate
\end{verbatim}

\textbf{Description}

Ensures the configured schema exists for the active database.

\subsection{novel backup create}

\textbf{Syntax}
\begin{verbatim}
novel backup create
\end{verbatim}

\textbf{Description}

Creates a JSON snapshot backup of the database contents.

\section{2. Novel Management}

\subsection{novel create}

\textbf{Syntax}
\begin{verbatim}
novel create "<novel>" [--genre <genre>] [--status <status>]
\end{verbatim}

\textbf{Description}

Creates a new novel entry.

\textbf{Example}
\begin{verbatim}
novel create "The Great Swordman" --genre fantasy --status planning
\end{verbatim}

\subsection{novel list}

\textbf{Syntax}
\begin{verbatim}
novel list
\end{verbatim}

\textbf{Description}

Displays all novels.

\subsection{novel show}

\textbf{Syntax}
\begin{verbatim}
novel show "<novel>"
\end{verbatim}

\textbf{Description}

Displays detailed information and record counts for a novel.

\subsection{novel update}

\textbf{Syntax}
\begin{verbatim}
novel update "<novel>" [--status <status>] [--summary "<text>"] [--genre <genre>]
\end{verbatim}

\textbf{Description}

Updates metadata for a novel.

\section{3. Sections}

\subsection{novel section add}

\textbf{Syntax}
\begin{verbatim}
novel section add "<novel>" <section>
\end{verbatim}

\textbf{Description}

Creates a logical grouping such as \texttt{plot} or \texttt{world}.

\section{4. Notes}

\subsection{novel note add}

\textbf{Syntax}
\begin{verbatim}
novel note add "<novel>" [--section <section>] [--title "<title>"] [--body "<text>"]
\end{verbatim}

\textbf{Description}

Creates a new note. When a section is provided and does not yet exist, it is created automatically.

\subsection{novel note append}

\textbf{Syntax}
\begin{verbatim}
novel note append <id> "<text>"
\end{verbatim}

\textbf{Description}

Appends content to an existing note.

\subsection{novel note tag}

\textbf{Syntax}
\begin{verbatim}
novel note tag <id> <tag>
\end{verbatim}

\textbf{Description}

Adds a tag to a note.

\section{5. Characters}

\subsection{novel character add}

\textbf{Syntax}
\begin{verbatim}
novel character add "<novel>" "<name>" [--role <role>]
\end{verbatim}

\textbf{Description}

Creates a character.

\subsection{novel character trait add}

\textbf{Syntax}
\begin{verbatim}
novel character trait add "<novel>" "<name>" <trait>
\end{verbatim}

\textbf{Description}

Adds a trait to a character.

\section{6. Relationships}

\subsection{novel relation add}

\textbf{Syntax}
\begin{verbatim}
novel relation add "<novel>" "<char1>" "<char2>" <type>
\end{verbatim}

\textbf{Description}

Creates a relationship between characters.

\section{7. Plot}

\subsection{novel plot add}

\textbf{Syntax}
\begin{verbatim}
novel plot add "<novel>" "<description>"
\end{verbatim}

\textbf{Description}

Adds a plot point.

\section{8. Chapters}

\subsection{novel chapter add}

\textbf{Syntax}
\begin{verbatim}
novel chapter add "<novel>" "<title>" [--number <num>]
\end{verbatim}

\textbf{Description}

Creates a chapter.

\section{9. Worldbuilding}

\subsection{novel world place add}

\textbf{Syntax}
\begin{verbatim}
novel world place add "<novel>" "<place>"
\end{verbatim}

\textbf{Description}

Adds a location.

\section{10. Timeline}

\subsection{novel timeline add}

\textbf{Syntax}
\begin{verbatim}
novel timeline add "<novel>" "<event>" [--day <num>]
\end{verbatim}

\textbf{Description}

Adds an event to the timeline.

\section{11. Search}

\subsection{novel search}

\textbf{Syntax}
\begin{verbatim}
novel search "<query>" [--type <entity>]
\end{verbatim}

\textbf{Description}

Searches across all stored data. Supported type filters include \texttt{novel}, \texttt{section}, \texttt{note}, \texttt{character}, \texttt{relation}, \texttt{plot}, \texttt{chapter}, \texttt{place}, \texttt{timeline}, and \texttt{quick}.

\section{12. Quick Capture}

\subsection{novel quick}

\textbf{Syntax}
\begin{verbatim}
novel quick "<novel>" "<text>"
\end{verbatim}

\textbf{Description}

Adds a fast, unstructured note.

\section{13. Reporting}

\subsection{novel review report}

\textbf{Syntax}
\begin{verbatim}
novel review report "<novel>"
\end{verbatim}

\textbf{Description}

Prints a compact inventory report for a novel.

\section{14. Example Workflow}

\begin{verbatim}
novel init
novel create "The Great Swordman"

novel character add "The Great Swordman" "John Ashen"
novel character add "The Great Swordman" "Jessica Vale"
novel relation add "The Great Swordman" "John Ashen" "Jessica Vale" rival

novel plot add "The Great Swordman" "John is expelled"
novel chapter add "The Great Swordman" "Chapter 1"

novel note add "The Great Swordman" --section world --title "Cultivation"
novel quick "The Great Swordman" "Hidden bloodline idea"

novel search "bloodline"
novel review report "The Great Swordman"
\end{verbatim}

\section{Conclusion}

The NKS CLI provides a scalable and structured approach to managing narrative data. It combines:

\begin{itemize}
    \item Structured SQL storage
    \item Graph-based character relationships
    \item Efficient CLI workflows
\end{itemize}

This enables both rapid idea capture and deeper structural organization.

\end{document}
